% Ported verbatim from paper/top3_joint_moment_bound.tex lines 214-442.
% Only cross-reference wording was adapted; no mathematics was altered.
We may work on any standard nonatomic probability space. All changes of
measure below can be realized on such a space.
\begin{lemma}\label{ranktwo}
Every $[0,1]$-valued graphon satisfies $\Sigma_2(W)\le7/6$.
\end{lemma}
\begin{proof}
We give the elementary argument for the rank-two absolute-energy
bound needed below.
Take any $\R^2$-valued law $\E YY^T=I_2$, write $Y=RU$ with a unit
choice at zero, $C=\E R$, and $N=\E[R UU^T]$.
Then $C^2\le2$ and
$\tr(N^2)\le1$, the latter from
$\E\|U-NY\|^2=1-\tr(N^2)\ge0$.
On the circle,
$g_{4,2}(t)=t^4-t^2+1/8$ is a positive kernel: if
$U=(\cos a,\sin a)$, it is $\cos(4(a-a'))/8$.
Hence its $RR'$-weighted integral is nonnegative.
The exact majorant
\[
 \frac29+\frac{11}{9}t^2-\frac49t^4-|t|
     =\frac{(1-|t|)(2|t|-1)^2(|t|+2)}9\ge0
\]
therefore gives
\[
 \E|Y\cdot Y'|\le\frac5{18}C^2+\frac79\tr(N^2)\le\frac43.
\]
Also $|\E Y|^2\le1$. For any orthonormal two-frame in a graphon,
$W(Y\cdot Y')\le(Y\cdot Y')_+$ pointwise, so its trace is at most
$(4/3+1)/2=7/6$. Taking the supremum proves the lemma.
The absolute-energy estimate is classical; see~\cite[Theorems 2 and 4]{Glazyrin}.
\end{proof}


\begin{lemma}\label{attainment}
The supremum $\mathcal M$ is attained, and
\[
 \cstar:=\frac{9+\sqrt5+2\sqrt{35}}{18}\le\mathcal M
       \le\frac{1+\sqrt3}{2}.
\]
Every maximizing graphon has three positive leading eigenvalues.
\end{lemma}
\begin{proof}
For an orthonormal $k$-frame, $K(x,y)=F(x)\cdot F(y)$ satisfies
$\int K=|\int F|^2\le1$ and $\int K^2=k$. Hence
$\int WK\le\int K_+\le(1+\sqrt k)/2$.
For the lower bound, use a six-part kernel with template $K_1\vee C_5$
and diagonal blocks equal to one. Give the center mass $a=1-5b$ and each
rim part mass $b$, where $0<b<1/5$. Put $\phi=(1+\sqrt5)/2$.
Its quotient spectrum is
\[
 \frac{1-2b\pm\sqrt{1+4b-36b^2}}2,
 \quad\phi b,\ \phi b,\ -b/\phi,\ -b/\phi.
\]
Indeed the four nonconstant Fourier modes on the pentagon give the last
four values; the remaining block is
$\left(\begin{smallmatrix}a&\sqrt{5ab}\\\sqrt{5ab}&3b\end{smallmatrix}\right)$.
Its positive eigenvalue exceeds $3b$, and its determinant is $-2ab<0$.
Consequently the sum of its three positive eigenvalues is
\[
 \frac{1+2\sqrt5 b+\sqrt{1+4b-36b^2}}2.
\]
Substitution of $b=(7+5\sqrt7)/126$ gives $\cstar$.
This is the usual quotient calculation for a join construction;
see~\cite{Hjoin}. In particular $\cstar>32/25>7/6$.

For completeness we record why $\Sigma_3$ of this step kernel is exactly the
sum of the three positive quotient eigenvalues. If $f$ is supported in one
part $A_i$ and $\int_{A_i}f=0$, then for $x\in A_i$ one has
$(T_Wf)(x)=\int_{A_i}f=0$, while for $x$ in any other part $A_j$ one has
$(T_Wf)(x)=W_{ji}\int_{A_i}f=0$; hence $T_Wf=0$. The whole block-mean-zero
subspace therefore lies in the zero eigenspace, and $\Sigma_3$ is governed
by the quotient matrix alone.

For attainment, let $D=W_1-W_2$. If $\lambda\ne0$ is an eigenvalue of
$T_D$ with unit eigenfunction $f$, then $|D|\le1$ gives
$\|f\|_\infty\le1/|\lambda|$. The cut-norm dual estimate yields
\[
 |\lambda|=|\langle f,T_Df\rangle|
 \le4\|D\|_\square\|f\|_\infty^2,
 \qquad \|T_D\|_{2\to2}^3\le4\|D\|_\square.
\]
The assertion is trivial when $T_D=0$; otherwise compactness supplies an
eigenvalue of modulus $\|T_D\|$. By the variational definition,
$|\Sigma_3(W_1)-\Sigma_3(W_2)|\le3\|T_D\|$. This also holds after
relabeling. Cut-metric compactness~\cite{LS} therefore gives a maximizer.
Finally Lemma~\ref{ranktwo} implies
$\lambda_3\ge\cstar-7/6>0$ at each maximizer.
\end{proof}

\begin{remark}\label{rem-newton}
Priority for the constant $\cstar$. The template $K_1\vee C_5$ of
Lemma~\ref{attainment} is the ``pole--pentagon'' law of the self-archived
Zenodo deposits~\cite{Newton1,Newton2}, which appeared about two and a half
months before the first version of the present paper. Writing the apex mass
as $p^*=(91-25\sqrt7)/126$ and each of the five rim masses as
$r^*=(7+5\sqrt7)/126$, so that $p^*+5r^*=1$, those deposits obtain the
identical value
\[
 \alpha_3=\frac12+\frac{\sqrt5}{18}+\frac{\sqrt{35}}9
        =\frac{9+\sqrt5+2\sqrt{35}}{18}=\cstar
\]
at the identical optimising weight, and they further assert
$\Lambda_3(G)\le\alpha_3n-3$ for every $n\ge3$. That assertion is
conditional. Its theorem is titled a \emph{candidate} universal $S_3$
theorem, it is stated to hold ``within the exact analytic and
proof-carrying certificate chain'', and the text declares that acceptance
``properly remains subject to independent mathematical audit of the
admissibility of the archived dual payments, the global sector partition and
the exact verifier implementation''. No refereed version is available. By
contrast, the computation in Lemma~\ref{attainment} is an ordinary quotient
calculation for a join and uses none of that machinery, so the lower bound
$\cstar\le\mathcal M$ stands independently; we cite~\cite{Newton1,Newton2}
for the constant and for the naming and claim no priority for either.
Nothing else in this paper depends on the unverified upper bound. Should it
later be confirmed unconditionally, then $\mathcal M=\cstar$ and the
interval displayed in Chapter~
ef{sec-further} collapses to a point.
\end{remark}

Fix a maximizing $W$ and a top-three orthonormal eigenframe $F$, taking
$f_1\ge0$. Write $L=\diag(\lambda_1,\lambda_2,\lambda_3)$.
Taking the absolute value of a top Rayleigh maximizer supplies the
nonnegative choice of $f_1$. We do not initially assume any eigenvalue simplicity.
\begin{lemma}\label{ellipse-structure}
Almost everywhere, $F^TLF=\mathcal M$. Moreover
$\lambda_3>\lambda_4$, where the fourth positive eigenvalue is zero if
absent. The canonical support $S=\supp(F_\#dx)$ is compact, lies on the
ellipsoid $v^TLv=\mathcal M$, and avoids the origin.
\end{lemma}
\begin{proof}
For a bounded real $g$ of mean zero, change density to $1+tg$ for both
signs of sufficiently small $t$. Set $Q=\int gFF^T$. The unchanged frame
has complete Gram matrix $I+tQ$ and numerator matrix
$L+t(QL+LQ)+O(t^2)$. After a common orthonormalization its trace is
\[
 \tr\bigl[(I+tQ)^{-1}\{L+t(QL+LQ)+O(t^2)\}\bigr]
   =\mathcal M+t\tr(LQ)+O(t^2).
\]
Maximality gives $\int gF^TLF=0$. Since every bounded mean-zero $g$ is
allowed, $F^TLF$ equals its mean $\mathcal M$ almost everywhere.
This applies to every active eigenframe. If
$\lambda_3=\lambda_4=\lambda>0$, compare the frames with third entries
$f_3$, $f_4$ and $(f_3+f_4)/\sqrt2$, keeping $f_1,f_2$ fixed.
The three identities imply $f_3^2=f_4^2$ and $f_3f_4=0$ almost everywhere,
a contradiction. Only finitely many null sets are combined.
Since $L\succ0$, the ellipsoid is compact and excludes zero. It contains
the full support, not just almost every frame value.
\end{proof}

\begin{lemma}\label{row-insertion}
For every measurable $h:[0,1]\to[0,1]$, the vector $r=\int hF$ satisfies
$r^TL^{-1}r\le\mathcal M$.
\end{lemma}
\begin{proof}
Give the original space $A$ mass $u=1-\epsilon$ and a new independent part
$B$ mass $\epsilon$; the kernel is $W$ on $A\times A$, equals the prescribed
cross-row $h$ (a function on $A$) on $A\times B$ and $B\times A$, and is zero
on $B\times B$. Put $t=L^{-1}r$ and $T=|t|^2$. For the trial frame take $f_i$
on $A$ and the constant $t_i$ on $B$; then the $A\times B$ and $B\times A$
contributions to the numerator are $\epsilon u\,r_it_j$ and its transpose.
Dividing the entire frame by $\sqrt u$ leaves every Rayleigh quotient
unchanged, and the resulting Gram and numerator matrices are
\[
 G=I+(\epsilon/u)tt^T,\qquad
 N=uL+\epsilon(rt^T+tr^T).
\]
The rank-one inverse formula gives the exact trial trace
\[
 \tr(G^{-1}N)=u\mathcal M+
       \frac{\epsilon r^TL^{-1}r}{1+\epsilon T/u}.
\]
This cannot exceed $\mathcal M$. Divide by $\epsilon$ and let
$\epsilon\downarrow0$.
\end{proof}

\begin{lemma}\label{contact-structure}
For $\nu=F_\#dx$ and every $v\in S$,
\begin{equation}\label{all-rows}
 \nu\{w:v\cdot w=0\}=0,\qquad
 \int_{v\cdot w>0}w\,d\nu(w)=Lv,\qquad v_1>0.
\end{equation}
The original kernel is $\mathbf1_{\{F(x)\cdot F(y)>0\}}$ almost
everywhere. It is irreducible, and $\lambda_1>\lambda_2$.
\end{lemma}
\begin{proof}
Let $Z=\{\int hF:0\le h\le1\}$ and
$E=\{r:r^TL^{-1}r\le\mathcal M\}$. The set $Z$ is compact by weak-star
compactness and is contained in $E$ by Lemma~\ref{row-insertion}.
Its support function is $H(v)=\int(v\cdot F)_+$.
For almost every $x$, the original row has integral vector $LF(x)$,
and $F(x)^TLF(x)=\mathcal M$. Since the support function of $E$ is
$\sqrt{\mathcal M v^TLv}$, it follows that $H(F(x))=\mathcal M$.
Continuity then gives $H(v)=\mathcal M$ for every $v\in S$.

For such $v$, the unique support point of the strictly convex ellipsoid
$E$ in direction $v$ is $Lv$. The row $h_v=\mathbf1_{\{v\cdot F>0\}}$
attains this support value, so $\int h_vF=Lv$. If
$A_0=\{v\cdot F=0\}$, then $h_v+\mathbf1_A$ has the same support value
for every measurable $A\subset A_0$. Uniqueness forces $\int_A F=0$
for every such $A$, hence $F=0$ almost everywhere on $A_0$.
The ellipsoid identity implies that $A_0$ is null.

It remains to identify every maximizing row, not merely those obtained by
perturbing $h_v$ on $A_0$. Let $h\in[0,1]$ satisfy $\int hF=Lv$ for some
$v\in S$. Pairing with $v$ and using $v^TLv=\mathcal M$ gives
$v\cdot\int hF=\mathcal M$, while on the other hand
\[
 v\cdot\int hF=\int h\,(v\cdot F)\,dx\le\int(v\cdot F)_+\,dx=H(v)=\mathcal M .
\]
Hence equality holds throughout. Since $0\le h\le1$, the pointwise bound
$h\,(v\cdot F)\le(v\cdot F)_+$ is therefore an equality almost everywhere,
which forces $h=1$ almost everywhere on $\{v\cdot F>0\}$ and $h=0$ almost
everywhere on $\{v\cdot F<0\}$. As $A_0$ is null, $h=h_v$ almost
everywhere. Thus the maximizing row is unique modulo null sets, which
identifies every original row and proves the threshold formula.

If a nontrivial measurable cut had zero cross-kernel, its indicator
multiplier would commute with $T_W$ and with the isolated top-three
spectral projection. The kernel $F(x)\cdot F(y)$ of that projection
would vanish across the cut, contradicting the zero-equator assertion
and Fubini. Thus $W$ is irreducible. The zero set of the nonnegative
$f_1$ would also give a zero cut by its eigen-equation, so $f_1>0$ almost
everywhere. If another top eigenfunction orthogonal to $f_1$ existed,
it would change sign. Equality in the absolute-value Rayleigh comparison
would force zero edges between its positive and negative sets; its
absolute value is a nonnegative top eigenfunction and is positive almost
everywhere. This contradicts irreducibility, proving simplicity.
Finally, for every $v\in S$, $\lambda_1v_1=\int h_vf_1>0$:
otherwise $h_v=0$ almost everywhere and hence $Lv=0$, which is impossible.
\end{proof}
